\begin{textsample}{NEG Dim 3 – global\_north\_2024\_10 – Score -1.00 – global\_north\_2024\_10/115977503.txt}  \label{ex:f3_neg_303}
Israel ' s war with Hezbollah in Lebanon is shattering young lives as conflict spreads in the Middle East Bekaa Valley, Lebanon -- The carnage of Israel ' s war with Hezbollah -- a conflict playing out in parallel and with direct links to the devastating war in the Gaza Strip -- continued over the weekend, with lives lost on both sides of the Israel-Lebanon border.
In Gaza, health officials said Monday that the toll from the war sparked by the Palestinian enclave ' s Hamas rulers, with their brutal Oct. 7, 2023 attack on Israel, had reached almost 42,300, with nearly 99,000 others wounded.
But while fighting in the decimated Gaza Strip continues, the Israeli military made a determined shift to what it calls the northern front in its broader war with Iran-backed groups in the region about a month ago.
Since then, Lebanese officials say more than 2,300 people have been killed in the country, and almost 10,700 others wounded.
The country ' s health ministry says 51 people were killed on Sunday alone.
Much of Israel ' s firepower has been directed @ @ @ @ @ @ @ @ @ @ and across the south of Lebanon.
Major Israeli ground operations in the south have also put United Nations \textbf{peacekeepers} in the line of fire.
But airstrikes have also pummelled Lebanon ' s eastern Bekaa Valley -- often with no warning.
Last week, CBS News visited the region ' s Rayak Hospital, which has been treating some of the youngest victims of the expanding war, including 16-year-old Ali Jaddouh.
Ali Jaddouh, 16, lays in a bed at the Rayak Hospital in Lebanon ' s Bekaa Valley, in early October 2024.
CBS News/Agnes Reau He recently had at least one badly damaged kidney and his colon at least partially removed, as well as his right leg above the knee.
He was in critical condition, with a dialysis machine doing the job of his shattered organs.
He told us he was in pain, and his haunted eyes suggested it was more than just physical.
The teenager said he was at home with his family late in the morning when an Israeli airstrike slammed into their @ @ @ @ @ @ @ @ @ @ have struck about 10 yards from where they were sitting."I wanted to run and help my mother, but I saw my leg was cut.
I lost consciousness and I don't recall what happened next,"he said.
He woke up in the hospital to find he ' d lost most of his leg."I was told that my father might be dead.
My mother can't walk anymore -- she lost her leg and had some damage in her back, and my eldest brother ' s face is burned."Israel says it is targeting Hezbollah in Lebanon, the Iran-backed, U.S. and Israeli-designated terrorist group that ' s fired thousands of rockets and drones at Israel since Oct. 8, 2023.
That ongoing barrage, said by Hezbollah to be in support of the Palestinian people and Iran ' s other allies, Hamas, includes the deadly drone attack over the weekend that killed four Israeli troops at a base in central Israel and wounded scores of other people.
Rayak Hospital @ @ @ @ @ @ @ @ @ @ weeks, the facility had treated only civilians.
While our team was there over the weekend, there was another Israeli strike nearby.
The agonising screams of two wounded girls who ' d been rushed onto the emergency ward echoed through the halls.
Nurse Mountaha Mkahal CBS News/Agnes Reau Nurse Mountaha Mkahal has been so busy looking after patients that, like many of the staff at Rayak, she ' s been sleeping at the hospital."It ' s very hard and unsettling,"she told CBS News."I am morally obliged to be here in war time -- not only to do my job when there is safety and peace.
This is the crucial time."She knows that with airstrikes coming often and without warning, members of her own young family could come through the emergency room doors any time.
Mkahal said seeing children suffer was the hardest part of her job.
Children like six-year-old Sawsan, who was brought in with six fractures to her skull.
Doctors had to remove @ @ @ @ @ @ @ @ @ @ so much pain that not even her mother ' s loving touch could ease the hurt or erase the horror.
Sawsan, 6, is seen in the Rayak Hospital in Lebanon ' s Bekaa Valley, where she was treated for multiple skull fractures and had shrapnel from an Israeli airstrike removed for her brain, in early October 2024.
CBS News/Agnes Reau"It ' s very hard to see a child suffering, and it reminds me of my own children, but I am hopefull those children and people will recover and go back to normal.
We have to do our best to help them recover, regardless of whether they can recover completely or not,"said the nurse.
The embassy noted that the country ' s primary airport in Beirut remained open and that flights out were still available commercially, and it noted that the U.S. government had helped hundreds of citizens leave on additional flights, warned those"additional flights will not continue indefinitely."It urged U.S. nationals in the country to complete and @ @ @ @ @ @ @ @ @ @ receive information about options to leave Lebanon."U.S. citizens who choose not to depart at this time should prepare contingency plans should the situation deteriorate further,"the embassy said, warning that those plans"should not rely on the U.S. government for assisted departure or evacuation."The U.S.
State Department has had a"Do Not Travel"advisory in effect for Lebanon for weeks, urging Americans not to visit the country in light of the Israel-Hezbollah violence.
Debora Patta is a CBS News foreign correspondent based in Johannesburg.
Since joining CBS News in 2013, she has reported on major stories across Africa, the Middle East and Europe.
Edward R.
Murrow and Scripps Howard awards are among the many accolades Patta has received for her work.

% matched lemmas: peacekeeper
\end{textsample}
